\documentclass[12pt,a4paper]{report}
\usepackage[T1]{fontenc}
\usepackage[utf8]{inputenc}
\usepackage{lmodern}
\usepackage[top=1cm, left=1cm, right=1cm]{geometry}
\usepackage[french]{babel}
\usepackage{listings}
\usepackage{hyperref}
\usepackage{graphicx}
\usepackage{fancyhdr}
\usepackage{lastpage}
\usepackage[most]{tcolorbox}
\usepackage{tikz}
\usepackage{amssymb}
\usepackage{enumitem}
\usepackage{exercise}



\usetikzlibrary{decorations.pathreplacing, backgrounds, calc}
\definecolor{QuizzColor}{RGB}{66, 209, 245}
\definecolor{AstuceColor}{RGB}{214, 212, 62}
\definecolor{DefinitionColor}{RGB}{0, 204, 102}
\definecolor{ComprendreColor}{RGB}{175, 122, 197}

\newtcolorbox{Quizz}{breakable,enhanced,colback=yellow!25!black!10!white,colframe={QuizzColor},drop fuzzy shadow,watermark color=white,title=Quizz - A vos buzzers ?,fonttitle=\large\bfseries,before skip=30pt plus 2pt,after skip=30pt plus 2pt,enlarge top by=0.25cm,%
overlay={\node[anchor=center, outer sep=-16pt] at (frame.north west) {\Huge}; }}

\newtcolorbox{Astuce}{breakable,enhanced,colback=yellow!25!black!10!white,colframe={AstuceColor},drop fuzzy shadow,watermark color=white,title=Astuce - Le saviez-vous ?,fonttitle=\large\bfseries,before skip=30pt plus 2pt,after skip=30pt plus 2pt,enlarge top by=0.25cm,%
overlay={\node[anchor=center, outer sep=-16pt] at (frame.north west) {\Huge}; }}

\newtcolorbox{Definition}[1][]{breakable,enhanced,colback=yellow!25!black!10!white,colframe={DefinitionColor},drop fuzzy shadow,watermark color=white,title=Définition - \textbf{{#1}},fonttitle=\large\bfseries,before skip=30pt plus 2pt,after skip=30pt plus 2pt,enlarge top by=0.25cm,%
overlay={\node[anchor=center, outer sep=-16pt] at (frame.north west) {\Huge}; }}

\newtcolorbox{Comprendre}{breakable,enhanced,colback=yellow!25!black!10!white,colframe={ComprendreColor},drop fuzzy shadow,watermark color=white,title=Vous avez dit Eureka ?,fonttitle=\large\bfseries,before skip=30pt plus 2pt,after skip=30pt plus 2pt,enlarge top by=0.25cm,%
overlay={\node[anchor=center, outer sep=-16pt] at (frame.north west) {\Huge}; }}


\newcommand{\motcle}[1]{\textbf{\textcolor[HTML]{ff8000}{#1}}}

\newcommand{\lstsansserif}{\ttfamily}
\lstdefinestyle{pythonstyle}{
    backgroundcolor=\color{backcolour},   
    commentstyle=\color{codegreen},
    keywordstyle=\color{magenta},
    numberstyle=\tiny\color{codegray},
    stringstyle=\color{codepurple},
    basicstyle=\large\lstsansserif,
    breakatwhitespace=false,         
    breaklines=true,                 
    captionpos=b,                    
    keepspaces=true,                 
    numbers=left,                    
    numbersep=5pt,                  
    showspaces=false,                
    showstringspaces=false,
    showtabs=false,                  
    tabsize=2,
}


% CSS
\lstdefinelanguage{CSS}{
  keywords={color,background-image, background,margin,padding,font,weight,display,position,top,left,right,bottom,list,style,border,size,white,space,min,width, transition:, transform:, transition-property, transition-duration, transition-timing-function,text-align},	
  sensitive=true,
  morecomment=[l]{//},
  morecomment=[s]{/*}{*/},
  morestring=[b]",
  alsoletter={:},
  alsodigit={-}
}

% JavaScript
\lstdefinelanguage{JavaScript}{
  morekeywords={typeof, new, true, false, catch, function, return, null, catch, switch, var, if, in, while, do, else, case, break},
  morecomment=[s]{/*}{*/},
  morecomment=[l]//,
  morestring=[b]",
}

% HTML
\lstdefinelanguage{HTML5}{
  language=html,
  sensitive=true,	
  alsoletter={<>=-},	
  morecomment=[s]{<!-}{-->},
  tag=[s],
  otherkeywords={
  % General
  >,
  % Standard tags
	<!DOCTYPE, <!doctype,
  </html, <html, <head, <title, </title, <style, </style, <link, </head, <meta, />,
	% body
	</body, <body,
	% Divs
	</div, <div, </div>, 
	% Paragraphs
	</p, <p, </p>, <h1, </h1, <h2, </h2, <h3, </h3, <h4, </h4,
	% scripts
	</script, <script,
  % More tags...
  <canvas, /canvas>, <svg, <rect, <animateTransform, </rect>, </svg>, <video, <source, <iframe, </iframe>, </video>, <image, </image>, <header, </header, <article, </article
  },
  ndkeywords={
  % General
  =,
  % HTML attributes
  charset=, src=, id=, width=, height=, style=, type=, rel=, href=,
  % SVG attributes
  fill=, attributeName=, begin=, dur=, from=, to=, poster=, controls=, x=, y=, repeatCount=, xlink:href=,
  % properties
  margin:, padding:, background-image:, border:, top:, left:, position:, width:, height:, margin-top:, margin-bottom:, font-size:, line-height:,color:, background:
	% CSS3 properties
  transform:, -moz-transform:, -webkit-transform:,
  animation:, -webkit-animation:,
  transition:,  transition-duration:, transition-property:, transition-timing-function:,
  }
}

\lstdefinestyle{htmlcssjsstyle} {%
  % General design
backgroundcolor=\color{backcolour},  
  basicstyle=\large\lstsansserif, 
  %frame=b,
  % line-numbers
  %xleftmargin={0.75cm},
  numbers=left,
  numberstyle=\tiny\color{codegray},
  stepnumber=1,
  firstnumber=1,
  numberfirstline=true,	
  % Code design
  identifierstyle=\color{black}\large\lstsansserif,
  keywordstyle=\color{blue},
  ndkeywordstyle=\color{codegreen},
  stringstyle=\color{editorOcher}\large\lstsansserif,
  commentstyle=\color{brown}\large\lstsansserif,
  % Code
  language=HTML5,
  alsolanguage=JavaScript,
  alsolanguage=CSS,
  alsodigit={.:;},	
  tabsize=2,
  showtabs=false,
  showspaces=false,
  showstringspaces=false,
  extendedchars=true,
  breaklines=true
}
%

\makeatother

% \lstset{style=htmlcssjsstyle, language=HTML5}
% \lstset{style=htmlcssjsstyle, language=JavaScript}
% \lstset{style=htmlcssjsstyle, language=CSS}
%\lstset{style=pythonstyle, language=Python}
\definecolor{codegreen}{rgb}{0,0.6,0}
\definecolor{codegray}{rgb}{0.5,0.5,0.5}
\definecolor{codepurple}{rgb}{0.58,0,0.82}
\definecolor{backcolour}{rgb}{0.95,0.95,0.92}
\definecolor{LightGray}{gray}{0.9}

\definecolor{lightgray}{rgb}{0.95, 0.95, 0.95}
\definecolor{darkgray}{rgb}{0.4, 0.4, 0.4}
%\definecolor{purple}{rgb}{0.65, 0.12, 0.82}
\definecolor{editorGray}{rgb}{0.95, 0.95, 0.95}
\definecolor{editorOcher}{rgb}{1, 0.5, 0} % #FF7F00 -> rgb(239, 169, 0)
\definecolor{editorGreen}{rgb}{0, 0.5, 0} % #007C00 -> rgb(0, 124, 0)
\definecolor{orange}{rgb}{1,0.45,0.13}		
\definecolor{olive}{rgb}{0.17,0.59,0.20}
\definecolor{brown}{rgb}{0.69,0.31,0.31}
\definecolor{purple}{rgb}{0.38,0.18,0.81}
\definecolor{lightblue}{rgb}{0.1,0.57,0.7}
\definecolor{lightred}{rgb}{1,0.4,0.5}
\renewcommand\listexercisename{Liste des exercices}%
\renewcommand\ExerciseName{Exercice}%
\renewcommand\AnswerName{Solution de l'exercice}%
\renewcommand\ExerciseListName{Ex.}%
\renewcommand\AnswerListName{Solution}%
\renewcommand\ExePartName{Partie}
\renewcommand\ExerciseSkipAfter{1cm}

\newcommand\emptybox[1]{%
 \hbox{\smash{\hspace*{\dimexpr-\hoffset}%
  }\framebox(\textwidth, #1){}}}


\newcommand*\touche[1]{%
  \begin{tikzpicture}[baseline=(key.base), very thin, line cap=round, black, rounded corners=0pt]%
    \node [draw, fill=white, fill opacity=1, rectangle, rounded corners=2pt, inner sep=1pt, minimum width=1.2em, font=\scriptsize\sffamily] (key) {#1\strut};

    \begin{scope}[on background layer]
      \draw [rounded corners=1pt, fill=white] ($ (key.north west) + (-2pt, 2pt) $) rectangle ($ (key.south east) + (2pt, -2pt) $);

      \fill [gray!60] ($ (key.south west) + (2pt, 0.1pt) $) -- ($ (key.south west) + (-1pt, -2pt) $)
                  -- ($ (key.south east) + (1pt, -2pt) $)  -- ($ (key.south east) + (-2pt, 0.1pt) $) -- cycle;

      \fill [gray!60] ($ (key.south east) + (-0.1pt, 2pt) $) -- ($ (key.south east) + (2pt, -1pt) $)
                  -- ($ (key.north east) + (2pt, 1pt) $)    -- ($ (key.north east) + (-0.1pt, -2pt) $) -- cycle;
    \end{scope}

    \draw ($ (key.north west) + (0.1pt, -2pt) $) -- ($ (key.north west) + (-2pt, 1pt) $);
    \draw ($ (key.north west) + (2pt, -0.1pt) $) -- ($ (key.north west) + (-1pt, 2pt) $);

    \draw ($ (key.north east) + (-0.1pt, -2pt) $) -- ($ (key.north east) + (2pt, 1pt) $);
    \draw ($ (key.north east) + (-2pt, -0.1pt) $) -- ($ (key.north east) + (1pt, 2pt) $);

    \draw ($ (key.south west) + (0.1pt, 2pt) $) -- ($ (key.south west) + (-2pt, -1pt) $);
    \draw ($ (key.south west) + (2pt, 0.1pt) $) -- ($ (key.south west) + (-1pt, -2pt) $);

    \draw ($ (key.south east) + (-0.1pt, 2pt) $) -- ($ (key.south east) + (2pt, -1pt) $);
    \draw ($ (key.south east) + (-2pt, 0.1pt) $) -- ($ (key.south east) + (1pt, -2pt) $);
  \end{tikzpicture}%
}

\usepackage{pgfplots}
\pgfplotsset{compat=1.18}

\title{Les algorithmes de recherche}

\date{}

\lstset{extendedchars=true,literate={é}{{\'e}}1}
\usepackage{fancyhdr}
\fancypagestyle{plain}{\fancyhf{}\fancyfoot[C]{\footnotesize Babash - NSI Tani Malandi et Johanne Bakalara \textperiodcentered{} CC BY-SA 4.0 \textperiodcentered{} page \thepage}\renewcommand{\headrulewidth}{0pt}}
\pagestyle{plain}
\RequirePackage{geometry}\geometry{a4paper,top=1.5cm,bottom=2.8cm,left=2cm,right=2cm,footskip=1.2cm}
\AtBeginDocument{\hypersetup{pdfauthor={Babash - NSI Tani Malandi et Johanne Bakalara},pdftitle={Complexité, tris et recherche : cours et TD-TP},pdfcreator={LuaLaTeX}}}
\begin{document}

\chapter*{ Mission : Optimisation Haute Fréquence}

\begin{Comprendre}
\textbf{CONTEXTE :}  Les serveurs de l'agence saturent ! Votre mission : déterminer quels algorithmes sont capables de supporter des flux massifs de données. Un \textbf{algorithme de recherche} sert à localiser une information dans une base de données.
    \\ \textbf{Protocole :} Vous allez analyser le code, mesurer le temps \textbf{ET} compter le nombre d'opérations élémentaires. 
\end{Comprendre}

\begin{lstlisting}[language=Python]
def recherche_lineaire(tab, valeur):
    for i in range(len(tab)):
        if tab[i] == valeur :
            return i # Cible trouvee !
    return None # Cible NON trouvee !
\end{lstlisting}

\begin{Exercise}[title={Check-up technique (Rappels de base) }]
    \Question Soit \texttt{serveurs = [102, 405, 201, 80, 199]}. Donnez le résultat de :
    \begin{itemize}
        \item \texttt{serveurs[1]} : \dotfill \quad \texttt{serveurs[-1]} : \dotfill \quad \texttt{serveurs[1:3]} : \dotfill
    \end{itemize}

    \Question On souhaite doubler la valeur du troisième élément de la liste \texttt{tab}. Quelle instruction écrire ? \dotfill

    \Question \textbf{L'accès direct :} Quel est l'impact sur le processeur d'un accès type \texttt{serveurs[0]} ? \textit{(Voir le cours)} \dotfill
    
    \Question \textbf{La recherche linéaire :} On parcourt la liste du début à la fin. Si la \texttt{valeur} recherchée se trouve à la dernière position d'une liste de taille $n$, combien d'étapes faut-il ? \dotfill
\end{Exercise}

\begin{Exercise}[title={Analyse du Code }]
    \Question Observez attentivement le code ci-dessous :
    \begin{itemize}
        \item \textbf{Ligne 7 (\texttt{milieu = ...}) :} Pourquoi utilise-t-on une division entière (\texttt{//}) ? \dotfill
        \item \textbf{Lignes 10 à 12 :} Pourquoi peut-on ignorer toute la partie droite si \texttt{tab[milieu] > valeur} ? \\
        .\dotfill\\
        .\dotfill
    \end{itemize}
\end{Exercise}


\begin{Comprendre}
    \textbf{Principe :} Si le tableau est \textbf{trié}, on regarde au milieu. Si la cible est plus grande, on ignore la moitié gauche. On divise ainsi la zone de recherche par deux à chaque étape.
    
\begin{lstlisting}[language=Python, numbers=left]
def recherche_dichotomique(tab, valeur):
    # Attention tab doit etre trié
    gauche = 0
    droite = len(tab) - 1
    while gauche <= droite:
        # On cible le centre de la zone de recherche
        milieu = (gauche + droite) // 2  
        if tab[milieu] == valeur: 
            return milieu # Cible trouvee !
        elif tab[milieu] > valeur:
            # Cible plus petite : on ignore la partie droite
            droite = milieu - 1
        else: 
            # Cible plus grande : on ignore la partie gauche
            gauche = milieu + 1
    return None # Cible NON trouvee !
\end{lstlisting}
\begin{center}
    \begin{tikzpicture}[x=1cm,y=1cm,c/.style={draw,minimum size=0.7cm,inner sep=0pt,font=\small},
 mn/.style={fill=orange!45},ok/.style={fill=gray!30},cl/.style={fill=cyan!30},ix/.style={thick,draw=black,line width=1.2pt},hors/.style={text=gray,draw=gray!60},lab/.style={anchor=west,font=\footnotesize,align=left}]
\node[lab] at (0,0.6){Recherche de $36$ dans \texttt{[3, 8, 11, 15, 19, 24, 30, 36, 41]}};
\node[font=\scriptsize,gray] at (0.00,0.1){0};
\node[font=\scriptsize,gray] at (0.70,0.1){1};
\node[font=\scriptsize,gray] at (1.40,0.1){2};
\node[font=\scriptsize,gray] at (2.10,0.1){3};
\node[font=\scriptsize,gray] at (2.80,0.1){4};
\node[font=\scriptsize,gray] at (3.50,0.1){5};
\node[font=\scriptsize,gray] at (4.20,0.1){6};
\node[font=\scriptsize,gray] at (4.90,0.1){7};
\node[font=\scriptsize,gray] at (5.60,0.1){8};
\node[c,] at (0.00,-0.60) {3};
\node[c,] at (0.70,-0.60) {8};
\node[c,] at (1.40,-0.60) {11};
\node[c,] at (2.10,-0.60) {15};
\node[c,mn] at (2.80,-0.60) {19};
\node[c,] at (3.50,-0.60) {24};
\node[c,] at (4.20,-0.60) {30};
\node[c,] at (4.90,-0.60) {36};
\node[c,] at (5.60,-0.60) {41};
\node[lab] at (6.6,-0.60){\texttt{gauche}=0, \texttt{droite}=8, \texttt{milieu}=4 : $19<36$,\\on garde la partie droite.};
\node[c,hors] at (0.00,-1.60) {3};
\node[c,hors] at (0.70,-1.60) {8};
\node[c,hors] at (1.40,-1.60) {11};
\node[c,hors] at (2.10,-1.60) {15};
\node[c,hors] at (2.80,-1.60) {19};
\node[c,] at (3.50,-1.60) {24};
\node[c,mn] at (4.20,-1.60) {30};
\node[c,] at (4.90,-1.60) {36};
\node[c,] at (5.60,-1.60) {41};
\node[lab] at (6.6,-1.60){\texttt{gauche}=5, \texttt{droite}=8, \texttt{milieu}=6 : $30<36$,\\on garde la partie droite.};
\node[c,hors] at (0.00,-2.60) {3};
\node[c,hors] at (0.70,-2.60) {8};
\node[c,hors] at (1.40,-2.60) {11};
\node[c,hors] at (2.10,-2.60) {15};
\node[c,hors] at (2.80,-2.60) {19};
\node[c,hors] at (3.50,-2.60) {24};
\node[c,hors] at (4.20,-2.60) {30};
\node[c,mn] at (4.90,-2.60) {36};
\node[c,] at (5.60,-2.60) {41};
\node[lab] at (6.6,-2.60){\texttt{gauche}=7, \texttt{droite}=8, \texttt{milieu}=7 : $36=36$,\\la cible est trouvée à l'indice 7.};
\node[c,mn] at (0,-3.9){};\node[lab] at (0.5,-3.9){case du milieu};\node[c,hors] at (3.5,-3.9){};\node[lab] at (4.0,-3.9){case écartée};
\end{tikzpicture}
\end{center}
\end{Comprendre}

\begin{Exercise}[title={Simulation de terrain}]
    On cherche la cible \textbf{17} dans : \texttt{L = [2, 5, 8, 12, 17, 21, 25]}.

    \Question \textbf{Méthode Linéaire :} Combien d'étapes pour trouver 17 ? \dotfill
    
    \begin{center}
    \begin{tabular}{|c|c|c|c|c|c|}
    \hline
    Étape & \texttt{gauche} & \texttt{droite} & \texttt{milieu} & \texttt{L[milieu]} & Test \\ \hline
    1 & 0 & 6 & 3 & 12 & $12 < 17$ (trop petit) \\ \hline
    2 & & & & & \\ \hline
    3 & & & & & \\ \hline
    \end{tabular}
    \end{center}

    \Question \textbf{Nombre de comparaisons :} Soit deux appels :
    \begin{itemize}
        \item \textit{recherche\_dichotomique([1,3,5,6,8,11,13,14,14,17,19,21,23], 19)}
        \item \textit{recherche\_lineaire([1,3,5,6,8,11,13,14,14,17,19,21,23], 19)}
    \end{itemize}
    Donner toutes les comparaisons faites dans chacun des cas.
    
    \Question \textbf{Analyse des cas :}
    \begin{itemize}
        \item \textbf{Meilleur Cas :}
        \begin{itemize}
            \item Pour la méthode linéaire, où doit être la cible pour finir en 1 étape ? \dotfill
            \item Pour la dichotomie, où doit être la cible pour finir en 1 étape ? \dotfill
        \end{itemize}
        \item \textbf{Pire Cas :}
        \begin{itemize}
            \item Quel est le pire des cas pour les deux algorithmes ? \dotfill
            \item Combien d'étapes max pour une liste de taille $100$ en linéaire ? \dotfill
            \item Combien d'étapes max pour une liste de taille $100$ en dichotomique ? \dotfill
            \item Combien d'étapes max pour une liste de taille $n$ en linéaire ? \dotfill
            \item Combien d'étapes max pour une liste de taille $n$ en dichotomique ? \dotfill
        \end{itemize}
    \end{itemize}
    Conclure sur la classe de complexité des deux algorithmes.
\end{Exercise}

\begin{Exercise}[title={Benchmark : De 10 à 100 Millions}]
    \Question \textbf{Expérimentation :} Lancez \texttt{recherche\_contest.py}. Notez les résultats (Pire cas) :
    \begin{center}
    \begin{tabular}{|l|c|c|c|c|}
    \hline
    Taille ($n$) & \multicolumn{2}{c|}{Recherche Linéaire } & \multicolumn{2}{c|}{Recherche Dichotomique } \\ \hline
    & Temps (s) & Opérations & Temps (s) & Opérations \\ \hline
    10 & & & & \\ \hline
    1 000 & & & & \\ \hline
    1 000 000 & & & & \\ \hline
    100 000 000 & & & & \\ \hline
    \end{tabular}
    \end{center}
    
    \Question \textbf{Observation :} Pour $n=10^8$, la dichotomie a nécessité environ \dotfill opérations.\\
    Est-ce cohérent avec la fonction $\log_2(n)$ ? \dotfill
\end{Exercise}

\begin{Exercise}[title={ \textbf{À RETENIR :}Synthèse de l'Agent}]
    \textit{Complétez le rapport ci-dessous en utilisant les termes : \textbf{trié}, \textbf{linéaire} $O(n)$, \textbf{logarithmique} $O(log(n))$, \textbf{milieu}, \textbf{1}, \textbf{$n$}, \textbf{dichotomique}.}
    
    L'efficacité d'une recherche dépend de l'algorithme choisi et de l'organisation des données :\\
    
    1. La \textbf{recherche linéaire} consiste à parcourir le tableau un par un. Dans le \textbf{meilleur des cas}, la cible est en première position et nécessite \_\_\_\_\_\_ opération.\\
    Dans le \textbf{pire des cas}, elle nécessite \_\_\_\_\_\_ opérations.\\
    Sa complexité est dite \textbf{\rule[-2pt]{2.5cm}{0.4pt}}.
    
    2. La \textbf{recherche \rule[-2pt]{2.5cm}{0.4pt}} est beaucoup plus rapide car elle divise la zone de recherche par deux à chaque étape en comparant la cible avec la valeur du \rule[-2pt]{2.5cm}{0.4pt}.
    
    3. Cependant, cette méthode ne fonctionne \textbf{que} si le tableau est préalablement \textbf{\rule[-2pt]{2.5cm}{0.4pt}}.\\
    Sa complexité est alors \textbf{\rule[-2pt]{2.5cm}{0.4pt}} : même pour 100 millions de données, moins de 30 opérations suffisent.
\end{Exercise}
\end{document}